%% =====================================================================
%% الفصل الثاني: الإطار القانوني في مدونة الشغل المغربية
%% =====================================================================

\chapter{الإطار القانوني في مدونة الشغل المغربية}
\label{chap:chapitre2}

%% ─── مقدمة الفصل ──────────────────────────────────────────
\section*{مقدمة الفصل}
\phantomsection\addcontentsline{toc}{section}{مقدمة الفصل}

أفرد المشرع المغربي في مدونة الشغل (القانون رقم \textenglish{65.99}) منظومة معيارية خاصة لتقنين العلاقات المهنية غير النمطية، مستهدفاً تحقيق التوازن بين متطلبات المقاولة في المرونة التسييرية والقدرة التنافسية من جهة، وبين صيانة النظام العام الاجتماعي وضمان الأمن المهني للأجراء من جهة ثانية. غير أن استقراء النصوص القانونية المنظمة لتقنيتي الشغل المؤقت والمناولة يكشف عن تفاوت هيكلي ملحوظ؛ إذ في الوقت الذي أحاط فيه المشرع الشغل المؤقت بترسانة تفصيلية من القيود الشكلية والموضوعية والمدد الزمنية الصارمة، اكتفى في المقابل بصياغة قواعد عامة ومقتضبة في تنظيم المقاولة من الباطن، مما فتح ثغرات للتحايل في الواقع الميداني.

وللإحاطة بهذا البناء التشريعي، يتناول هذا الفصل أربعة محاور مترابطة:
\begin{itemize}[leftmargin=*, itemsep=4pt]
    \item \textbf{أولاً:} المرجعيات القانونية وهيكلتها في الدستور والمعايير الدولية والمدونة؛
    \item \textbf{ثانياً:} أحكام التشغيل المؤقت في مدونة الشغل (المواد \textenglish{495} إلى \textenglish{506})؛
    \item \textbf{ثالثاً:} أحكام المناولة وعقود الشغل المحددة المدة (المواد \textenglish{86}--\textenglish{91} و\textenglish{16}--\textenglish{17})؛
    \item \textbf{رابعاً:} النصوص التنظيمية وقراءة نقدية في مكامن القصور التشريعي.
\end{itemize}

%% ---------------------------------------------------------------------
\section{المرجعيات القانونية}

\subsection{المرجعيات الدستورية والمعايير الدولية}
تتأصل القواعد الحاكمة للتشغيل غير المباشر في المنظومة الدستورية للمملكة المغربية لسنة \textenglish{2011}؛ حيث ينص الفصل \textenglish{31} على التزام الدولة والمؤسسات العمومية والجماعات الترابية بتعبئة كل الوسائل المتاحة لتيسير أسباب استفادة المواطنات والمواطنين على قدم المساواة من الحق في الشغل والحماية الاجتماعية، فيما يكفل الفصل \textenglish{35} حرية المبادرة والمقاولة والتنافس الحر.

وتتكامل هذه المرجعية الدستورية الوطنية مع معايير منظمة العمل الدولية (\textenglish{OIT})، وعلى رأسها \textbf{الاتفاقية رقم \textenglish{181} بشأن وكالات الاستخدام الخاصة (\textenglish{1997})} التي تُلزم بتقييد الوساطة وإخضاعها للترخيص وحظر اقتطاع أي عمولات من أجور العمال، و\textbf{التوصية رقم \textenglish{198} بشأن علاقة العمل (\textenglish{2006})} التي تكرس مبدأ «أولوية الواقع» عند تكييف العقود الشغلية.

\subsection{التوزيع البنيوي للنصوص داخل مدونة الشغل}
اعتمد المشرع المغربي منهجية الفصل الهيكلي بين مؤسستي الشغل المؤقت والمناولة داخل مدونة الشغل (القانون رقم \textenglish{65.99} الصادر بتنفيذه الظهير الشريف رقم \textenglish{1.03.194} بتاريخ \textenglish{11} أغسطس \textenglish{2003}):
\begin{itemize}[leftmargin=*, itemsep=5pt]
    \item \textbf{أحكام المناولة (المقاولة من الباطن):} أُدرجت ضمن الكتاب الأول المتعلق بالاتفاقيات المتعلقة بالشغل (المواد \textenglish{86} إلى \textenglish{91})، بوصفها عقداً تجارياً ومهنياً يكلف فيه مقاول أصلي مقاولاً من الباطن بإنجاز شغل أو خدمة.
    \item \textbf{أحكام الشغل المؤقت:} أُفردت لها المواد \textenglish{495} إلى \textenglish{506} ضمن \textbf{الباب الثاني من الكتاب الرابع} المتعلق بـ «الوساطة في الاستخدام وتشغيل الأجراء»، مما يعكس فلسفة المشرع في التعامل معها كنشاط استثنائي للوساطة خاضع لأحكام التراخيص والمراقبة الإدارية المقررة في المواد \textenglish{477} وما يليها.
\end{itemize}

%% ---------------------------------------------------------------------
\section{أحكام التشغيل المؤقت (المواد 495 وما يليها)}

\subsection{شروط ممارسة النشاط ونظام الترخيص والكفالة المالية}
انطلاقاً من الطبيعة الحساسة لنشاط التشغيل المؤقت وما يكتنفه من مخاطر المضاربة بقوة العمل، أخضع المشرع تأسيس وتسيير هذه المقاولات لضوابط مالية وإدارية صارمة:
\begin{itemize}[leftmargin=*, itemsep=5pt]
    \item \textbf{الترخيص الإداري المسبق والرأسمال الأدنى (المادتان \textenglish{477} و\textenglish{481}):} لا يجوز فتح مقاولة للتشغيل المؤقت إلا بعد الحصول على إذن مسبق من السلطة الحكومية المكلفة بالشغل، ويشترط ألا يقل رأسمالها عن \textenglish{100.000} درهم مع انتفاء الأحكام الماسة بالشرف أو الأمانة في حق المسيرين.
    \item \textbf{إيداع الكفالة المالية بصندوق الإيداع والتدبير (المادتان \textenglish{482} و\textenglish{488}):} ألزم المشرع في المادة \textenglish{482} كل مقاولة تشغيل مؤقت بإيداع كفالة لدى صندوق الإيداع والتدبير (\textenglish{CDG}) تعادل \textenglish{50} مرة القيمة الإجمالية السنوية للحد الأدنى القانوني للأجر؛ وتسمح المادة \textenglish{488} للمحكمة المختصة، عند إعسار الوكالة أو سحب ترخيصها، بتخصيص هذه الكفالة لأداء مستحقات الأجراء وواجبات الصندوق الوطني للضمان الاجتماعي (\textenglish{CNSS}).
    \item \textbf{سحب الترخيص واستمرار الالتزامات (المادتان \textenglish{487} و\textenglish{503}):} يجوز للسلطة الحكومية سحب الإذن بالممارسة بقرار معلل ودون تعويض (المادة \textenglish{487})، غير أن السحب لا يُعفي مسؤولي المقاولة من الوفاء بالتزاماتهم تجاه الأجراء و\textenglish{CNSS} (المادة \textenglish{503}).
\end{itemize}

\subsection{الحالات الحصرية والمدد القصوى والوقاية والسلامة (المواد 496--505)}
رسخت المادة \textenglish{496} مبدأً جوهرياً يقضي بأن اللجوء إلى الشغل المؤقت هو \textbf{استثناء محصور في أربع حالات فقط} بعد استشارة الهيئات التمثيلية للأجراء: إحلال أجير محل آخر (في غير حالة الإضراب)، أو التزايد المؤقت في النشاط، أو الأشغال الموسمية، أو الأشغال التي استقر العرف على عدم اللجوء فيها لعقد غير محدد المدة.

كما حظرت المادة \textenglish{497} اللجوء إلى الأجراء المؤقتين لإنجاز أشغال تكتسي خطورة خاصة، ومنعت المادة \textenglish{498} استخدامهم لمواجهة تزايد مؤقت في المناصب التي شملها فصل لأسباب اقتصادية خلال السنة الموالية للفصل. وألزمت المادة \textenglish{504} المقاولة المستعملة باتخاذ كافة تدابير الصحة والسلامة والتأمين عن حوادث الشغل والأمراض المهنية (وفق القانون رقم \textenglish{18.12}). ويوضح الجدول الآتي ضوابط المدة وفترة التجربة والعقوبة الزجرية:

\begin{table}[htbp]
\centering
\small
\renewcommand{\arraystretch}{1.55}
\begin{tabularx}{\textwidth}{|p{4.5cm}|p{4.5cm}|X|}
\hline
\textbf{حالة اللجوء (المادة \textenglish{496})} & \textbf{المدة القصوى للمهمة (المادة \textenglish{500})} & \textbf{فترة التجربة (المادة \textenglish{502})} \\
\hline
\textbf{إحلال أجير محل آخر} & طيلة مدة توقف عقد الأجير المتغيب & يومان إذا كان العقد مبرماً لأقل من شهر \\
\cline{1-2}
\textbf{تزايد مؤقت في نشاط المقاولة} & \textenglish{3} أشهر قابلة للتجديد مرة واحدة فقط & \textenglish{3} أيام إذا تراوحت مدة العقد بين شهر وشهرين \\
\cline{1-2}
\textbf{أشغال موسمية وأعراف مهنية} & \textenglish{6} أشهر غير قابلة للتجديد & \textenglish{5} أيام إذا فاقت مدة العقد شهرين \\
\hline
\multicolumn{3}{|r|}{\textbf{الجزاء الجنائي (المادة \textenglish{505}):} غرامة مالية من \textenglish{2.000} إلى \textenglish{5.000} درهم عن مخالفة أحكام هذا الباب.} \\
\hline
\end{tabularx}
\caption{الضوابط الإجرائية والمدد الزمنية وفترة التجربة في الشغل المؤقت}
\label{tab:durees-temporaire}
\end{table}

%% ---------------------------------------------------------------------
\section{أحكام المناولة وعقود الشغل المحددة المدة}

\subsection{أحكام عقد المقاولة من الباطن (المواد 86 إلى 91)}
نصت المادة \textenglish{87} من مدونة الشغل على أن المقاول من الباطن يُعتبر مشغلاً خاضعاً لكافة أحكام المدونة، ونصوص الضمان الاجتماعي، والتعويض عن حوادث الشغل والأمراض المهنية. وأوجبت المادة \textenglish{88} تعليق اسم المقاول الأصلي وعنوانه في كل ورش أو منشأة وتحديد بيانات الأداء وبطاقة الشغل وفق الشكليات القانونية متى اشتغل الأجراء في منشآت المقاول الأصلي.

أما بشأن \textbf{المسؤولية الاحتياطية للمقاول الأصلي}، فقد نصت المادتان \textenglish{89} و\textenglish{90} على أنه إذا كان المقاول من الباطن:
\begin{itemize}[leftmargin=*, itemsep=4pt]
    \item \textbf{غير مقيد في السجل التجاري}؛
    \item \textbf{ولا يتوفر على أصل تجاري}؛
    \item \textbf{وكان في حالة إعسار}؛
\end{itemize}
أصبح المقاول الأصلي مسؤولاً—في حدود المبالغ المترتبة بذمته لفائدة المقاول من الباطن—عن أداء الأجور، والتعويض عن العطلة السنوية، وتعويضات الفصل، واشتراكات \textenglish{CNSS}، والرسم الخاص بالتكوين المهني، مع تخويل الأجراء و\textenglish{CNSS} حق إقامة دعوى مباشرة ضده (المادة \textenglish{90}). كما يسأل المقاول الأصلي بصفة مباشرة عن التعويض عن حوادث الشغل والأمراض المهنية إذا أنجزت الأشغال أو الخدمات داخل مؤسساته أو ملحقاتها (الفقرة الأخيرة من المادة \textenglish{89}).

\subsection{ضوابط عقد الشغل المحدد المدة (المادتان 16 و17)}
ينظم القسم الخاص بانعقاد الشغل عقد الشغل المحدد المدة (\textenglish{CDD}) كآلية مرونة ثنائية مباشرة بين المشغل والأجير دون وسيط:
\begin{itemize}[leftmargin=*, itemsep=4pt]
    \item \textbf{شروط الصحة (المادة \textenglish{16}):} لا يُبرم عقد محدد المدة إلا في حالات استثنائية (إحلال أجير، ازدياد مؤقت في النشاط، عمل موسمي)، أو في القطاعات والحالات الاستثنائية المحددة بنص تنظيمي، أو عند فتح مقاولة لأول مرة أو إطلاق منتوج جديد (لمدة أقصاها سنة قابلة للتجديد مرة واحدة في القطاعات غير الفلاحية).
    \item \textbf{التحويل التلقائي (المادة \textenglish{17}):} ينص القانون صراحة على أن عقد الشغل المحدد المدة يتحول بقوة القانون إلى عقد غير محدد المدة إذا استمر الأجير في عمله بعد انتهاء المدة المحددة فيه، وهي قاعدة صريحة في الـ \textenglish{CDD} لكنها غير منصوص عليها بنص مماثل في باب التشغيل المؤقت.
\end{itemize}

%% ---------------------------------------------------------------------
\section{النصوص التنظيمية وقراءة نقدية في مكامن القصور}

\subsection{المراسيم التطبيقية والاتفاقيات الجماعية}
تعزز الإطار التشريعي للمدونة بمراسيم تطبيقية، إلى جانب دور الاتفاقيات الجماعية (المواد \textenglish{104} وما يليها) في تحسين ظروف الشغل القطاعية:
\begin{enumerate}[leftmargin=*, itemsep=5pt]
    \item \textbf{المرسوم رقم \textenglish{2.19.456} (الصادر في \textenglish{30} يناير \textenglish{2020}، الجريدة الرسمية عدد \textenglish{6855}):} المحدد لتكوين وطريقة عمل اللجنة المختصة بمقاولات التشغيل المؤقت المنصوص عليها في المادة \textenglish{496}، والتي تتألف من \textenglish{18} عضواً وفق تركيبة ثلاثية (الإدارة، المنظمات المهنية للمشغلين، والمنظمات النقابية للأجراء الأكثر تمثيلاً)، وقد حل محل المرسوم السابق رقم \textenglish{2.04.464} الصادر سنة \textenglish{2004}.
    \item \textbf{المرسوم رقم \textenglish{2.19.793} (الصادر في \textenglish{23} يوليو \textenglish{2020}):} المحدد للقطاعات والحالات الاستثنائية التي يمكن فيها إبرام عقد الشغل محدد المدة استناداً إلى المادة \textenglish{16} من المدونة.
\end{enumerate}

\subsection{مكامن القصور التشريعي والثغرات الهيكلية}
تكشف القراءة النقدية لأحكام مدونة الشغل عن ثغرات بنيوية تتكرر آثارها في الممارسة الميدانية:
\begin{itemize}[leftmargin=*, itemsep=5pt]
    \item \textbf{اختلال التوازن بين الشغل المؤقت والمناولة:} مقابل التقييد الصارم للشغل المؤقت بمدد لا تتعدى \textenglish{3} إلى \textenglish{6} أشهر (المادة \textenglish{500})، وردت أحكام المناولة (المواد \textenglish{86}--\textenglish{91}) فضفاضة دون تحديد مدد قصوى أو حالات حصرية، مما شجع على التحايل عبر «مناولة مقنعة».
    \item \textbf{ضيق شروط المادة \textenglish{89}:} يعفي اشتراط عدم تقييد مقاول الباطن في السجل التجاري وعدم امتلاكه أصلاً تجارياً المقاولَ الأصلي من المسؤولية بمجرد إنشاء شركة شكلية مقيدة بالسجل التجاري، كما يتيح حصر الأداء «في حدود ما تبقى في الذمة» التحايل بتصفية الفواتير مسبقاً.
    \item \textbf{غياب المساواة الأجرية الصريحة وضعف العقوبات:} تخلو المدونة من نص صريح يُلزم بالمساواة في الأجر بين الأجير المؤقت والأجير الدائم (عكس القانون الفرنسي)، كما أن غرامة المادة \textenglish{505} (\textenglish{2.000} إلى \textenglish{5.000} درهم) ضعيفة ولا توازيها قاعدة تشريعية صريحة تقرر إعادة التكييف التلقائي.
\end{itemize}

%% ─── خلاصة الفصل ──────────────────────────────────────────
\section*{خلاصة الفصل}
\phantomsection\addcontentsline{toc}{section}{خلاصة الفصل}

يتضح مما سبق أن الإطار القانوني المغربي يضع قواعد دقيقة للشغل المؤقت وأحكاماً عامة للمناولة، غير أنه يشكو من ثغرات تشريعية أبرزها ضيق نطاق المادة \textenglish{89}، وهزالة غرامة المادة \textenglish{505}، وسكوت النص عن المساواة الأجرية وعن إعادة التكييف التلقائي. وللوقوف على الضوابط التي قيد بها المشرع اللجوء إلى هذه الآليات، ننتقل في الفصل الثالث إلى تحليل الحالات الحصرية التي يسمح فيها القانون باللجوء إلى الشغل المؤقت.
